\section{A proof architecture for micro-to-macro claims}
\label{sec:book-micro-macro-architecture}

This appendix instantiates the six obligations in
Eq.~\ref{eq:book-cross-scale-proof-obligations}.  It serves as a proof benchmark; its hard-sphere
variables carry no proposed identification with protein variables.

A statement that microscopic motion produces a macroscopic law contains several distinct claims.
The microscopic dynamics must be specified on a complete state space.  The macroscopic observable
must be a declared projection of that state.  A scaling limit must distinguish leading histories
from negligible ones.  The leading sector must close under an autonomous equation, and every
discarded sector must carry a quantitative remainder bound.  Irreversibility requires one further
account of where information leaves the retained description.

The long-time derivation of the Boltzmann equation from dilute hard-sphere dynamics by Deng, Hani,
and Ma supplies a particularly complete benchmark \cite{denghani2025boltzmann}.  Its value here is
methodological.  The proof turns a broad emergence claim into a sequence of exact interfaces:
correlation marginals, connected cumulants, collision-history graphs, local geometric integrals,
topological complexity bounds, and closure onto one kinetic equation.

\subsection{The theorem fixes the observable, scaling, and error}

Consider hard spheres of diameter $\varepsilon$ in dimension $d\geq2$.  The theorem uses the
specific grand-canonical initial density whose $N$-particle sector is proportional to
\begin{equation}
\varepsilon^{-(d-1)N}
\prod_{j=1}^{N}f_0(z_j)\,
\mathbf1_{\mathcal D_N^\varepsilon}(z_N),
\label{eq:book-hard-sphere-initial-ensemble}
\end{equation}
where $f_0\geq0$ and $\int f_0=1$.  This choice gives the Boltzmann--Grad scaling
\begin{equation}
\mathbb E N\,\varepsilon^{d-1}\asymp1.
\label{eq:book-boltzmann-grad-scaling}
\end{equation}
Let $f_s^\varepsilon(t,z_s)$ denote the rescaled $s$-particle correlation function and let
$\mathcal D_s^\varepsilon$ be the hard-core admissible set.  Fix $\beta>0$.  Suppose the hard-sphere
Boltzmann equation has a solution $f(t,x,v)$ on $[0,t_{\mathrm{fin}}]$ satisfying
\begin{equation}
\sup_{0\leq t\leq t_{\mathrm{fin}}}
\|e^{2\beta|v|^2}f(t)\|_{L^\infty_{x,v}}\leq A,
\qquad
\|f_0\|_{\mathrm{Bol}_{2\beta}}
+\|\nabla_xf_0\|_{\mathrm{Bol}_{2\beta}}\leq B_0.
\label{eq:book-hard-sphere-solution-assumptions}
\end{equation}
Then, for $\varepsilon$ sufficiently small depending on
$(d,t_{\mathrm{fin}},\beta,A,B_0)$,
\begin{equation}
\sup_{\substack{0\leq t\leq t_{\mathrm{fin}}\\
                  1\leq s\leq |\log\varepsilon|}}
\left\|
f_s^\varepsilon(t,z_s)
-\mathbf1_{\mathcal D_s^\varepsilon}(z_s)
 \prod_{j=1}^{s} f(t,z_j)
\right\|_{L^1}
\leq \varepsilon^\theta ,
\label{eq:book-hard-sphere-long-time-chaos}
\end{equation}
where $\theta>0$ depends only on dimension.

Equation~\ref{eq:book-hard-sphere-long-time-chaos} is a propagation-of-chaos statement for low-order
statistics.  It compares a specific microscopic marginal with a specific solution of a nonlinear
kinetic equation in a specified norm.  The time interval may be any fixed interval contained in
the regular lifespan of the Boltzmann solution.  With a globally regular and uniformly bounded
solution, the proof also permits $t_{\mathrm{fin}}$ to grow slowly as $\varepsilon\to0$.  The
theorem leaves global regularity of the Boltzmann equation and trajectory-wise convergence outside
its scope.

The phrase ``long time'' is relative to the convergence radius of the classical collision-tree
expansion.  A direct Duhamel expansion of either BBGKY or Boltzmann histories produces a series with
growth of the form $\sum_n C^n t^n$.  Matching the two series term by term therefore yields only a
short-time theorem.  Extending the result requires a representation that preserves the history of
correlations without expanding the leading factorized contribution through the entire interval.

\subsection{Partial expansion carries structure across time layers}

Split $[0,t_{\mathrm{fin}}]$ into short layers of length $\tau$.  At a layer boundary the truncated
correlation hierarchy has the schematic cumulant expansion
\begin{equation}
\widetilde f_s
=\sum_{H\subseteq[s]}
  (f^A)^{\otimes([s]\setminus H)}E_H
+\operatorname{Err},
\qquad E_\varnothing=1.
\label{eq:book-hard-sphere-cumulant-expansion}
\end{equation}
The one-particle term $f^A$ collects histories whose observed roots remain disconnected.  The
cumulant $E_H$ collects histories in which the roots indexed by $H$ become connected through
collisions, overlaps introduced by the expansion, or initial links.

The partial time expansion treats these two sectors asymmetrically.  On each layer, $f^A$ is
expanded only to the preceding boundary and is then compared with the Boltzmann solution.  The
older cumulants are expanded through preceding layers until their histories reach the initial
ensemble.  The leading series never accumulates the full $t^n$ proliferation.  The connected terms
retain the complete information needed to estimate their rarity and arrive with powers of
$\varepsilon$ that can pay for a longer expansion.

This construction also locates the time direction.  A layer-by-layer estimate using only a natural
Banach norm of $E_H$ cannot supply the induction: velocity reversal at an intermediate time leaves
such norms unchanged and lets the Newtonian dynamics run toward either temporal direction.  The
same hypothetical induction would then identify the initial marginal with a tensor product built
from a later Boltzmann state.  The proof instead transports the structured collision history back
to the distinguished, nearly factorized ensemble at time zero.  The initial class and the retained
history carry the orientation.

\subsection{Molecules turn histories into local probability estimates}

Each expanded history is reduced to a diagram called a molecule.  Atoms represent collisions or
overlaps, colored particle lines join successive events of one particle, and layers retain the
time interval containing each event.  The paper uses two related ranks.  Write
$\rho_{\mathrm{circ}}(M)$ for the circuit rank of one molecule.  For a history with layers
$M_{\ell'}$, define the bookkeeping complexity
\begin{equation}
\rho_{\mathrm{hist}}
=\sum_{\ell'=0}^{\ell}s_{\ell'}
+\sum_{\ell'=1}^{\ell}R_{\ell'},
\qquad
R_{\ell'}=\rho_{\mathrm{circ}}(M_{\ell'}),
\label{eq:book-molecule-history-rank}
\end{equation}
where $s_{\ell'}$ counts particle lines carrying connected information into layer $\ell'$ and
$s_0$ counts lines entering through initial links.  These ranks are comparable up to the constants
used in the proof, while they serve different bookkeeping roles.  For fixed size and complexity,
the number of admissible molecules obeys a bound of the form
\begin{equation}
\#\mathcal M(|M|,\rho_{\mathrm{hist}})
\leq C^{|M|}|\log\varepsilon|^{C\rho_{\mathrm{hist}}}.
\label{eq:book-molecule-count}
\end{equation}

The nonnegative contribution attached to one diagram is then rewritten exactly as a local
integral.  For a full molecule with root set $H$,
\begin{equation}
\bigl\||\operatorname{IN}_M|\bigr\|_{L^1}
=\varepsilon^{(d-1)|H|} I_M(Q_M).
\label{eq:book-molecule-local-integral}
\end{equation}
Every atom contributes Dirac constraints imposing free transport, contact geometry, and elastic
scattering.  The factor $Q_M$ carries the admissible time domain, including layer membership and
event ordering, together with the input densities and support restrictions.  Locality matters:
each constraint sees the variables adjacent to one atom instead of a long chain of substituted
trajectories.

Cutting a molecule is a diagrammatic choice of Fubini order,
\begin{equation}
I_M=I_{M_1}\circ I_{M_2}.
\label{eq:book-molecule-fubini}
\end{equation}
Repeated cutting and splitting reduce the integral to elementary one- and two-atom pieces.  A
normal collision is neutral in powers of $\varepsilon$.  A bare recollision can also be neutral.
The robust gain unit is typically a nondegenerate collision followed by a recollision, denoted a
good $\{33\}$ piece: closing onto a previous trajectory forces a collision time, relative velocity,
or spatial displacement into a thin set and yields a factor $\varepsilon^\nu$.  Strong and weak
resonances are first split into separate subcases; their smallness can then be recovered from other
local $\{3\}$, $\{4\}$, $\{33\}$, or $\{44\}$ structures.  Certain $\{4\}$ pieces also carry a
loss $\varepsilon^{-(d-1)}$ because an integration lacks a fixed input.

For the main cumulant family, the central combinatorial theorem constructs an integration order
for which the global exponent has a positive balance,
\begin{equation}
\frac{\nu}{2}\#\mathrm{good}
+(d-1)\bigl(|H|-\#\mathrm{bad}\bigr)
\geq c\rho_{\mathrm{hist}}.
\label{eq:book-molecule-gain-balance}
\end{equation}
For the expansion-error family, the right side has the weaker form
$c\rho_{\mathrm{hist}}-C|H|$; the large number of atoms in those histories supplies a short-layer
time factor that absorbs the loss.  The estimates are global accounts rather than one-to-one
assignments of one gain to every recollision.  Combining
Eqs.~\ref{eq:book-molecule-count}--\ref{eq:book-molecule-gain-balance} turns historical complexity
into a net smallness factor of order $\varepsilon^{c\rho_{\mathrm{hist}}}$ in the diagram
contribution, with room to absorb the polylogarithmic number of diagrams.  This is the quantitative
step that makes the connected cumulants vanish.

\subsection{The surviving trees close onto the Boltzmann equation}

After the cyclic terms have been controlled, the leading one-particle contribution consists of
trees with $\rho_{\mathrm{circ}}(M)=0$.  The cluster expansion contains additional
collision/overlap-labelled trees.
An involution toggles the highest eligible collision or overlap atom, preserves its integral, and
reverses its sign.  Those terms cancel in pairs.  The remaining rooted binary trees are indexed by
the same recursion as the gain and loss terms in the Boltzmann Duhamel expansion.

The microscopic tree still collides at a displaced surface $x+\varepsilon\omega$, whereas the
Boltzmann iterate uses a point collision.  Spatial regularity controls that displacement.  Large
trees are suppressed by the short-layer time factor; cyclic trees are suppressed by the cutting
estimate.  Finally, the truncated dynamics is compared with the original dynamics.  An absolute
upper bound on the number of collisions among a fixed number of hard spheres forces sufficiently
many particles to participate when the recollision count is large, creating enough good local
pieces to make the discarded event rare.

The proof therefore has three closure estimates:
\begin{equation}
\boxed{
\begin{aligned}
f^A&=f_{\mathrm{Boltzmann}}+o(1),\\
E_H&=o(1)\quad(H\ne\varnothing),\\
f_s^{\mathrm{err}}&=o(1).
\end{aligned}}
\label{eq:book-three-boltzmann-closures}
\end{equation}
The first identifies the autonomous macroscopic law.  The second proves propagation of chaos.
The third reconnects the controlled expansion to the stated microscopic dynamics.

\subsection{Three information quantities must remain separate}

The finite hard-sphere flow is reversible and preserves its fine-grained Liouville information.
The complexity $\rho_{\mathrm{hist}}$ measures connected collision-history topology and appears in
a diagram estimate.  The limiting Boltzmann density has the $H$-functional
\begin{equation}
H_{\mathrm B}(f)=\int f\log f\,dx\,dv,
\label{eq:book-boltzmann-h-functional}
\end{equation}
whose monotonicity follows from the kinetic equation.  These three quantities live on different
state spaces and obey different laws.

The paper's velocity-reversal argument establishes why a norm-only induction cannot create this
direction.  A time-reversed evolved ensemble generally lies outside the special nearly factorized
initial class in Eq.~\ref{eq:book-hard-sphere-initial-ensemble}, and the theorem makes no convergence
claim for that reversed ensemble.  The direction enters through the initial data class and the
partial expansion anchored at time zero.  The retained one-particle limit then obeys an
entropy-directed kinetic semigroup while every finite microscopic trajectory remains reversible.

The diagram estimate nevertheless has a useful energy-like coordinate.  After the root prefactor,
local integrations, and subcase sum have been combined, write $\mathcal W_\rho$ for the resulting
nonnegative contribution at history complexity $\rho=\rho_{\mathrm{hist}}$.  Its
$\rho$-dependent bound has the schematic form
\begin{equation}
\mathcal W_\rho
\lesssim |\log\varepsilon|^{C\rho}\varepsilon^{c\rho}
=\exp\!\left[-c\rho\log(1/\varepsilon)
              +C\rho\log|\log\varepsilon|\right].
\label{eq:book-history-effective-action}
\end{equation}
Thus $-\log\mathcal W_\rho$ acts as an effective cost for diagram contributions.  The weights arise
inside an upper-bound expansion and do not form a normalized probability law over histories.  A
Gibbs law would require a separately defined ensemble, normalization, and matching lower bounds.

\subsection{What transfers from the benchmark}

The proof succeeds by keeping six interfaces explicit: microscopic state, retained observable,
scaling parameter, history expansion, quantitative remainder, and closed limiting law.  Its account
of irreversibility also identifies the information discarded by the retained marginal and the
special initial class that orients the limit.

Protein applications use the same proof discipline at different state spaces.  A fitted sequence
Hamiltonian needs a declared evolutionary sampling law; a residue interaction corrected for
phylogeny needs a generative row-tree null; a spectral mode needs an identified operator and a
perturbation bound; a structural or functional interpretation needs a matched endpoint.  Chapters
\ref{chap:book-phylogeny}, \ref{chap:book-observations}, \ref{chap:book-spectra}, and
\ref{chap:book-falsification} supply those domain-specific objects.  The hard-sphere result supplies
the standard of closure: every change of state space carries an explicit approximation and an error
whose scale can be checked.
